% ==============================================================================
% CHAPITRE 2 : ANALYSE ET CONCEPTION
% ==============================================================================
\chapter{Analyse et Conception}

\begin{infobox}[Objectif du chapitre]
Ce chapitre présente l'analyse fonctionnelle et non fonctionnelle du système,
les diagrammes UML, l'architecture globale, le modèle de données et la
justification des choix technologiques.
\end{infobox}

% ──────────────────────────────────────────────────────────────────────────────
\section{Analyse des Besoins}
% ──────────────────────────────────────────────────────────────────────────────

\subsection{Identification des acteurs}

Le système interagit avec deux catégories d'acteurs principaux :

\begin{table}[H]
\centering
\caption{Acteurs du système}
\label{tab:acteurs}
\renewcommand{\arraystretch}{1.4}
\begin{tabular}{>{\bfseries}p{3cm} p{5cm} p{5cm}}
\toprule
\rowcolor{bleuprincipal!15}
\textbf{Acteur} & \textbf{Rôle} & \textbf{Interactions principales} \\
\midrule
Étudiant & Soumetteur de commentaires & Soumettre un commentaire, consulter le résultat, voir l'historique \\
\midrule
\rowcolor{gray!8}
Enseignant & Consommateur d'analyses & Consulter le dashboard, filtrer par matière, voir les statistiques \\
\midrule
Système IA & Acteur interne & Analyser le sentiment, retourner les scores \\
\bottomrule
\end{tabular}
\end{table}

\subsection{Besoins fonctionnels}

Les besoins fonctionnels représentent les fonctionnalités que le système
doit obligatoirement offrir :

\begin{enumerate}[label=\textbf{BF\arabic*}., leftmargin=3em]
  \item Le système doit permettre à un étudiant de soumettre un commentaire
        textuel d'au moins 5 caractères.
  \item Le système doit analyser automatiquement le sentiment du commentaire
        et retourner la classe prédite (Négatif/Neutre/Positif) avec les
        scores de confiance.
  \item Le système doit permettre d'associer un commentaire à une matière
        académique spécifique.
  \item Le système doit sauvegarder chaque commentaire analysé avec ses
        métadonnées dans la base de données.
  \item Le système doit permettre la consultation de l'historique des
        commentaires avec pagination.
  \item Le système doit offrir des filtres de recherche par sentiment,
        par matière et par date.
  \item Le système doit permettre la suppression d'un commentaire par
        son identifiant.
  \item Le système doit exposer des statistiques agrégées sur les analyses
        effectuées.
  \item L'application mobile doit permettre la navigation entre l'écran
        d'analyse et l'écran historique.
  \item Le tableau de bord doit afficher des visualisations interactives
        des tendances de sentiments.
\end{enumerate}

\subsection{Besoins non fonctionnels}

\begin{table}[H]
\centering
\caption{Besoins non fonctionnels}
\label{tab:bnf}
\renewcommand{\arraystretch}{1.4}
\begin{tabular}{>{\bfseries}p{3cm} p{10cm}}
\toprule
\rowcolor{bleuprincipal!15}
\textbf{Critère} & \textbf{Exigence} \\
\midrule
Performance & Le temps de réponse d'une analyse ne doit pas dépasser 2 secondes en conditions normales \\
\midrule
\rowcolor{gray!8}
Disponibilité & L'API doit être accessible 24h/24 une fois déployée \\
\midrule
Scalabilité & L'architecture doit permettre d'ajouter de nouvelles matières sans modification du code \\
\midrule
\rowcolor{gray!8}
Maintenabilité & Le code doit être modulaire, documenté et versionné via Git \\
\midrule
Sécurité & Les données étudiantes doivent être stockées localement, sans transfert vers des services tiers \\
\midrule
\rowcolor{gray!8}
Compatibilité & L'application mobile doit fonctionner sur Android 8.0+ \\
\midrule
Portabilité & Le système doit être déployable sous Linux Ubuntu sans dépendances propriétaires \\
\bottomrule
\end{tabular}
\end{table}

% ──────────────────────────────────────────────────────────────────────────────
\section{Diagrammes UML}
% ──────────────────────────────────────────────────────────────────────────────

\subsection{Diagramme de cas d'utilisation}

\begin{figure}[H]
\centering
\begin{tikzpicture}[
  actor/.style={circle, draw=bleuprincipal, fill=bleuprincipal!10,
                minimum size=1cm, font=\small\bfseries},
  usecase/.style={ellipse, draw=bleuprincipal, fill=white,
                  minimum width=3.5cm, minimum height=0.8cm,
                  font=\footnotesize, align=center, thick},
  system/.style={rectangle, draw=bleuprincipal!60, dashed,
                 minimum width=10cm, minimum height=8cm,
                 label={[font=\bfseries\small, text=bleuprincipal]
                        above:{Système d'Analyse de Sentiments}}}
]
  % Acteurs
  \node[actor] (etudiant) at (-5, 0) {};
  \node[below=0.1cm of etudiant, font=\small\bfseries] {Étudiant};

  \node[actor] (enseignant) at (7, 0) {};
  \node[below=0.1cm of enseignant, font=\small\bfseries] {Enseignant};

  % Frontière système
  \node[system] (sys) at (1, 0) {};

  % Cas d'utilisation
  \node[usecase] (uc1) at (1,  3.0) {Soumettre un\\commentaire};
  \node[usecase] (uc2) at (1,  1.8) {Voir le résultat\\d'analyse};
  \node[usecase] (uc3) at (1,  0.6) {Consulter\\l'historique};
  \node[usecase] (uc4) at (1, -0.6) {Filtrer par\\sentiment/matière};
  \node[usecase] (uc5) at (1, -1.8) {Supprimer un\\commentaire};
  \node[usecase] (uc6) at (1, -3.0) {Voir les\\statistiques};

  % Relations Étudiant
  \draw[->] (etudiant) -- (uc1);
  \draw[->] (etudiant) -- (uc2);
  \draw[->] (etudiant) -- (uc3);
  \draw[->] (etudiant) -- (uc4);
  \draw[->] (etudiant) -- (uc5);

  % Relations Enseignant
  \draw[->] (enseignant) -- (uc3);
  \draw[->] (enseignant) -- (uc4);
  \draw[->] (enseignant) -- (uc6);
  \draw[->] (enseignant) -- (uc5);

\end{tikzpicture}
\caption{Diagramme de cas d'utilisation du système}
\label{fig:uml_usecase}
\end{figure}

\subsection{Diagramme de séquence — Analyse d'un commentaire}

\begin{figure}[H]
\centering
\begin{tikzpicture}[
  lifeline/.style={draw=bleuprincipal!50, dashed, thick},
  box/.style={rectangle, draw=bleuprincipal, fill=bleuprincipal!10,
              minimum width=2.2cm, minimum height=0.7cm,
              font=\small\bfseries, align=center},
  msg/.style={->, thick, color=bleuprincipal!80},
  ret/.style={->, thick, dashed, color=vertprincipal!80}
]
  % En-têtes
  \node[box] (user)    at (0,0)   {Étudiant\\(Flutter)};
  \node[box] (api)     at (4,0)   {FastAPI\\Backend};
  \node[box] (service) at (8,0)   {Service\\Prédiction};
  \node[box] (model)   at (11.5,0){CamemBERT\\Model};
  \node[box] (db)      at (15,0)  {PostgreSQL};

  % Lignes de vie
  \draw[lifeline] (0,-0.35) -- (0,-10);
  \draw[lifeline] (4,-0.35) -- (4,-10);
  \draw[lifeline] (8,-0.35) -- (8,-10);
  \draw[lifeline] (11.5,-0.35) -- (11.5,-10);
  \draw[lifeline] (15,-0.35) -- (15,-10);

  % Messages
  \draw[msg] (0,-1.2) -- node[above,font=\footnotesize]{POST /api/predict} (4,-1.2);
  \draw[msg] (4,-2.0) -- node[above,font=\footnotesize]{predict(texte)} (8,-2.0);
  \draw[msg] (8,-2.8) -- node[above,font=\footnotesize]{nettoyer(texte)} (8,-2.8);
  \draw[msg] (8,-3.5) -- node[above,font=\footnotesize]{tokeniser(texte)} (11.5,-3.5);
  \draw[msg] (11.5,-4.3) -- node[above,font=\footnotesize]{forward()} (11.5,-4.3);
  \draw[ret] (11.5,-5.1) -- node[above,font=\footnotesize]{logits} (8,-5.1);
  \draw[msg] (8,-5.9) -- node[above,font=\footnotesize]{softmax(logits)} (8,-5.9);
  \draw[ret] (8,-6.7) -- node[above,font=\footnotesize]{\{sentiment, scores\}} (4,-6.7);
  \draw[msg] (4,-7.5) -- node[above,font=\footnotesize]{INSERT commentaire} (15,-7.5);
  \draw[ret] (15,-8.3) -- node[above,font=\footnotesize]{id, date\_creation} (4,-8.3);
  \draw[ret] (4,-9.1) -- node[above,font=\footnotesize]{JSON résultat complet} (0,-9.1);

\end{tikzpicture}
\caption{Diagramme de séquence : analyse d'un commentaire}
\label{fig:seq_analyse}
\end{figure}

\subsection{Diagramme de classes}

\begin{figure}[H]
\centering
\begin{tikzpicture}[
  class/.style={rectangle, draw=bleuprincipal, fill=white,
                minimum width=5cm, thick},
  attr/.style={font=\small\ttfamily},
  meth/.style={font=\small\ttfamily, color=bleuprincipal}
]

% Classe Commentaire (DB)
\node[class] (commentaire) at (0,0) {
  \begin{tabular}{l}
    \rowcolor{bleuprincipal!20}
    \multicolumn{1}{c}{\textbf{Commentaire (DB)}} \\
    \hline
    \small\ttfamily id : Integer (PK) \\
    \small\ttfamily texte\_original : Text \\
    \small\ttfamily texte\_nettoye : Text \\
    \small\ttfamily matiere : String \\
    \small\ttfamily sentiment\_predit : Integer \\
    \small\ttfamily label\_sentiment : String \\
    \small\ttfamily score\_negatif : Float \\
    \small\ttfamily score\_neutre : Float \\
    \small\ttfamily score\_positif : Float \\
    \small\ttfamily date\_creation : DateTime \\
    \hline
    \small\ttfamily to\_dict() : dict \\
  \end{tabular}
};

% Classe ServicePrediction
\node[class] (service) at (8.5, 2) {
  \begin{tabular}{l}
    \rowcolor{vertprincipal!20}
    \multicolumn{1}{c}{\textbf{ServicePrediction}} \\
    \hline
    \small\ttfamily \_instance : Singleton \\
    \small\ttfamily modele : CamemBERT \\
    \small\ttfamily tokeniseur : Tokenizer \\
    \small\ttfamily device : torch.device \\
    \hline
    \small\ttfamily charger() : void \\
    \small\ttfamily predict(texte) : dict \\
  \end{tabular}
};

% Classe Commentaire (Flutter)
\node[class] (flutter) at (8.5, -2.5) {
  \begin{tabular}{l}
    \rowcolor{orangeaccent!15}
    \multicolumn{1}{c}{\textbf{Commentaire (Flutter)}} \\
    \hline
    \small\ttfamily id : int \\
    \small\ttfamily texteOriginal : String \\
    \small\ttfamily matiere : String? \\
    \small\ttfamily sentimentPredit : int \\
    \small\ttfamily scores : Map \\
    \hline
    \small\ttfamily fromJson(json) \\
    \small\ttfamily emoji : String \\
    \small\ttfamily scoreConfiance : double \\
  \end{tabular}
};

% Flèches
\draw[->, thick, color=bleuprincipal] (service) -- node[above, font=\small]{utilise} (commentaire);
\draw[->, thick, dashed, color=orangeaccent] (flutter) -- node[right, font=\small]{mappe} (commentaire);

\end{tikzpicture}
\caption{Diagramme de classes simplifié}
\label{fig:class_diagram}
\end{figure}

% ──────────────────────────────────────────────────────────────────────────────
\section{Architecture Globale du Système}
% ──────────────────────────────────────────────────────────────────────────────

\subsection{Vue d'ensemble}

\begin{figure}[H]
\centering
\begin{tikzpicture}[
  layer/.style={rectangle, rounded corners=6pt, minimum width=3.5cm,
                minimum height=1.2cm, align=center, font=\small\bfseries,
                thick, draw},
  arrow/.style={->, thick, color=bleuprincipal!70},
  label/.style={font=\tiny, color=gray}
]

% Couche présentation
\node[layer, fill=bleuclaire!20, draw=bleuclaire] (flutter) at (0, 6)
  {\faIcon{mobile-alt} Application\\Flutter};
\node[layer, fill=bleuclaire!20, draw=bleuclaire] (streamlit) at (5, 6)
  {\faIcon{chart-bar} Dashboard\\Streamlit};

% Couche API
\node[layer, fill=bleuprincipal!15, draw=bleuprincipal, minimum width=9cm]
  (fastapi) at (2.5, 4) {\faIcon{server} API REST — FastAPI (port 8000)};

% Couche services
\node[layer, fill=vertprincipal!15, draw=vertprincipal] (predict) at (0, 2)
  {\faIcon{brain} Service\\Prédiction};
\node[layer, fill=vertprincipal!15, draw=vertprincipal] (router) at (5, 2)
  {\faIcon{route} Routeur\\Commentaires};

% Couche IA
\node[layer, fill=orangeaccent!15, draw=orangeaccent] (camembert) at (0, 0)
  {\faIcon{robot} CamemBERT\\Fine-tuné};

% Couche données
\node[layer, fill=rougeprincipal!10, draw=rougeprincipal] (postgres) at (5, 0)
  {\faIcon{database} PostgreSQL\\sentiment\_etudiants};

% Connexions
\draw[arrow] (flutter) -- node[label, right]{HTTP/JSON} (fastapi);
\draw[arrow] (streamlit) -- node[label, left]{HTTP/JSON} (fastapi);
\draw[arrow] (fastapi) -- (predict);
\draw[arrow] (fastapi) -- (router);
\draw[arrow] (predict) -- (camembert);
\draw[arrow] (router) -- node[label, right]{SQLAlchemy ORM} (postgres);
\draw[arrow] (predict) -- node[label, below]{sauvegarder} (router);

% Légende
\node[font=\tiny, gray] at (2.5, -1.2)
  {Architecture 4 couches : Présentation — API — Métier — Données};

\end{tikzpicture}
\caption{Architecture globale du système en 4 couches}
\label{fig:architecture}
\end{figure}

\subsection{Description des couches}

\subsubsection{Couche Présentation}
Comprend deux interfaces utilisateur distinctes selon le profil d'utilisateur :
l'application mobile Flutter pour les étudiants, et le dashboard Streamlit
pour les enseignants/administrateurs. Les deux communiquent exclusivement
via l'API REST.

\subsubsection{Couche API (FastAPI)}
Le point d'entrée unique du système. FastAPI gère la validation des
requêtes entrantes via les schémas Pydantic, l'authentification CORS,
le routage vers les services appropriés et la sérialisation des réponses
JSON.

\subsubsection{Couche Métier}
Contient la logique fonctionnelle : le \texttt{ServicePrediction} (singleton
gérant le modèle CamemBERT) et le routeur de commentaires gérant les
opérations CRUD.

\subsubsection{Couche Données}
PostgreSQL stocke toutes les analyses effectuées. SQLAlchemy ORM fait
office de couche d'abstraction, permettant d'interagir avec la base via
des objets Python sans écrire de SQL brut.

% ──────────────────────────────────────────────────────────────────────────────
\section{Modèle de Données}
% ──────────────────────────────────────────────────────────────────────────────

\subsection{Modèle Conceptuel de Données (MCD)}

Le modèle de données est intentionnellement simple : une seule entité
principale \textbf{Commentaire} regroupe toutes les informations relatives
à une analyse de sentiment.

\begin{figure}[H]
\centering
\begin{tikzpicture}[
  entity/.style={rectangle, draw=bleuprincipal, fill=bleuprincipal!10,
                 minimum width=6cm, font=\small\bfseries},
  attr/.style={font=\footnotesize\ttfamily}
]
\node[entity] (header) at (0,0) {COMMENTAIRE};
\node[rectangle, draw=bleuprincipal!60, fill=white, minimum width=6cm,
      align=left, font=\footnotesize\ttfamily] (body) at (0,-3.5) {
  \begin{tabular}{ll}
    \textbf{PK} & id : SERIAL \\
    & texte\_original : TEXT \\
    & texte\_nettoye : TEXT \\
    & matiere : VARCHAR(100) \\
    & sentiment\_predit : INTEGER \\
    & label\_sentiment : VARCHAR(20) \\
    & score\_negatif : FLOAT \\
    & score\_neutre : FLOAT \\
    & score\_positif : FLOAT \\
    & date\_creation : TIMESTAMP \\
  \end{tabular}
};
\draw[bleuprincipal] (header.south west) -- (body.north west);
\draw[bleuprincipal] (header.south east) -- (body.north east);
\end{tikzpicture}
\caption{Modèle Conceptuel de Données}
\label{fig:mcd}
\end{figure}

\subsection{Modèle Logique de Données (MLD)}

Le MLD correspond directement à la table SQL créée par SQLAlchemy :

\begin{lstlisting}[style=styleSQL, caption={Schéma SQL de la table commentaires}, label={lst:sql_schema}]
CREATE TABLE commentaires (
    id               SERIAL          PRIMARY KEY,
    texte_original   TEXT            NOT NULL,
    texte_nettoye    TEXT,
    matiere          VARCHAR(100),
    sentiment_predit INTEGER         NOT NULL,  -- 0, 1 ou 2
    label_sentiment  VARCHAR(20)     NOT NULL,  -- 'Negatif', 'Neutre', 'Positif'
    score_negatif    DOUBLE PRECISION,
    score_neutre     DOUBLE PRECISION,
    score_positif    DOUBLE PRECISION,
    date_creation    TIMESTAMP       DEFAULT NOW()
);
\end{lstlisting}

% ──────────────────────────────────────────────────────────────────────────────
\section{Justification des Choix Technologiques}
% ──────────────────────────────────────────────────────────────────────────────

\begin{table}[H]
\centering
\caption{Justification des choix technologiques}
\label{tab:choix_tech}
\renewcommand{\arraystretch}{1.5}
\begin{tabular}{p{2.5cm} p{3cm} p{7cm}}
\toprule
\rowcolor{bleuprincipal!15}
\textbf{Technologie} & \textbf{Rôle} & \textbf{Justification} \\
\midrule
CamemBERT & Modèle NLP & État de l'art pour le français, pré-entraîné sur 138 Go, surpasse tous les modèles multilingues \\
\midrule
\rowcolor{gray!8}
PyTorch & Framework DL & Standard industriel, support natif HuggingFace, flexible et débuggable \\
\midrule
HuggingFace & Bibliothèque & Simplifie le chargement et fine-tuning des Transformers, vaste communauté \\
\midrule
\rowcolor{gray!8}
FastAPI & API REST & Le plus rapide des frameworks Python API, typage fort, documentation automatique Swagger \\
\midrule
SQLAlchemy & ORM & Abstraction DB portable, sécurité SQL injection, migrations faciles \\
\midrule
\rowcolor{gray!8}
PostgreSQL & Base de données & ACID, performances, support JSON natif, open-source et robuste \\
\midrule
Streamlit & Dashboard & Python pur, déploiement instantané, intégration Plotly native \\
\midrule
\rowcolor{gray!8}
Flutter & Application mobile & Cross-platform (Android/iOS/Web), Material 3, performances natives \\
\midrule
Dart & Langage mobile & Typage fort, compilation AOT, excellent pour les UIs réactives \\
\bottomrule
\end{tabular}
\end{table}

\section*{Conclusion du Chapitre}

L'analyse fonctionnelle a permis d'identifier précisément les besoins des
utilisateurs et de concevoir une architecture répondant à ces besoins de
manière optimale. Le choix de CamemBERT comme modèle central, combiné à
FastAPI, PostgreSQL et Flutter, forme un écosystème cohérent et moderne.
Le chapitre suivant détaille l'implémentation concrète de chaque composant.

\cleardoublepage
